\subsection{\texorpdfstring{\(L^{\infty}(G)\) 的弱闭不变子空间}{L\^{}\{\textbackslash infinity\}(G) 的弱闭不变子空间}}

在本节中，我们把 \(\mathrm{L}^{\infty}(\mathrm{G})\) 与 \(\mathrm{L}^{1}(\mathrm{G})\) 的对偶等同，并给它赋以弱拓扑 \(\sigma(\mathrm{L}^{\infty}(\mathrm{G}),\mathrm{L}^{1}(\mathrm{G}))\)。我们用 \((f,g)\mapsto\langle f,g\rangle\)（对 \(f\in\mathrm{L}^{1}(\mathrm{G})\) 与 \(g\in\mathrm{L}^{\infty}(\mathrm{G})\)）表示定义这一对偶的双线性映射。

映射 W \(\mapsto\) W° 是从 L∞(G) 的弱闭向量子空间所成的集合到 L¹(G) 的闭向量子空间所成的集合上的一个双射（EVT，II，第 55 页，命题 10）。

另一方面，若 \(f \in \mathrm{L}^1(\mathrm{G})\) 且 \(x \in \mathrm{G}\)，则自同态 \(g \mapsto f * g\)（相应地，\(g \mapsto \varepsilon_x * g\)）（它是 Banach 空间 \(\mathrm{L}^1(\mathrm{G})\) 的自同态）以自同态 \(h \mapsto \check{f} * h\)（相应地，\(h \mapsto \varepsilon_{x^{-1}} * h\)）（Banach 空间 \(\mathrm{L}^\infty(\mathrm{G})\) 的自同态）为其转置（INT，VIII，§4，第 3 节，例 6）。\(\mathrm{L}^1(\mathrm{G})\) 的一个闭向量子空间是 \(\mathrm{L}^1(\mathrm{G})\) 的理想，其必要充分条件是它在 G 的平移下不变。因此，\(\mathrm{L}^\infty(\mathrm{G})\) 的一个弱闭向量子空间在与 \(\mathrm{L}^1(\mathrm{G})\) 中元素的卷积下稳定，其必要充分条件是它在 G 的平移下不变。

设 W 是 \(L^{\infty}(G)\) 的一个弱闭向量子空间。假设 W（从而 \(W^{\circ}\) 也）在 G 的平移下不变。设 \(f \in L^{1}(G)\)。对每个 \(g \in L^{\infty}(G)\)，有 \((\check{f} * g)(x) = \langle \varepsilon_{x} * f, g \rangle = \langle f, \varepsilon_{x^{-1}} * g \rangle\)。因此，f 属于 \(W^{\circ}\) 的必要充分条件是 \(\check{f} * g = 0\) 对所有 \(g \in W\) 成立。

若 W 是 \(L^{\infty}(G)\) 的一个在 G 的平移下不变的弱闭向量子空间，我们用 A(W) 表示属于 W 的特征标 \(\chi \in \widehat{G}\) 所成的集合。这是 \(\widehat{G}\) 的一个闭子集。若 F 是 \(\widehat{G}\) 的一个闭子集，我们用 Y(F) 表示由 F 的元素生成的 \(L^{\infty}(G)\) 的弱闭向量子空间；由于 G 的每个平移把每个特征标变为与之成比例的函数，空间 Y(F) 是平移不变的。

设 W 是 \(\mathrm{L}^{\infty}(\mathrm{G})\) 的一个在 G 的平移下不变的弱闭子空间。由双极定理（EVT，II，第 48 页，定理 1），特征标 \(\chi\) 属于 W 当且仅当它属于 \((\mathrm{W}^{\circ})^{\circ}\)；后一空间是由这样的函数 \(g \in \mathrm{L}^{\infty}(\mathrm{G})\) 组成的集合：\(\langle f, g \rangle = 0\) 对每个 \(f \in W^{\circ}\) 成立。我们有 \(\langle f, \chi \rangle = \overline{\mathcal{F}}(f)(\chi)\)，从而

\[
\mathrm{A} (\mathrm{W}) = \mathrm{V} (\mathrm{W} ^ {\circ}).
\]

类似地，函数 \(f \in \mathrm{L}^1(\mathrm{G})\) 属于 \(\Upsilon(\mathrm{F})^\circ\) 当且仅当 \(\langle f, \chi \rangle = 0\) 对所有 \(\chi \in \mathrm{F}\) 成立，这等价于 \(\overline{\mathcal{F}}(f)(\chi) = 0\)（对 \(\chi \in \mathrm{F}\)），也就是 \(f \in \Upsilon(\mathrm{F})\)。因此（同前）我们有

\[
\mathrm{Y} (\mathrm{F}) = \Upsilon (\mathrm{F}) ^ {\circ}.
\]

关系式 \(V(\Upsilon(F)) = F\)（I，第 13 页及 I，第 30 页）与 \(\Upsilon(V(I)) \supset I\)，结合双极定理（EVT，II，第 48 页，定理 1），便蕴含

\[
\mathrm{A} (\mathrm{Y} (\mathrm{F})) = \mathrm{F}, \quad \mathrm{Y} (\mathrm{A} (\mathrm{W})) \subset \mathrm{W}.
\]

命题 4.　设 W 是 L∞(G) 的一个非零的弱闭平移不变向量子空间。则 W 至少含有 G 的一个特征标。

我们已看到 \(A(W) = V(W^{\circ})\)。由于 \(W \neq 0\)，有 \(W^{\circ} \neq L^{1}(G)\)，于是由 II，第 252 页定理 1，\(V(W^{\circ})\) 非空。

命题 5.　设 W 是 L∞(G) 的一个弱闭平移不变向量子空间。

\begin{enumerate}
\def\labelenumi{\alph{enumi})}
\item
  对 A(W) 在 \(\widehat{G}\) 中的任意邻域 U，W 中每个函数都是属于 U 的特征标的线性组合的一个弱极限；
\item
  若 A(W) 的边界不含任何非空完全集，则 W 的每个函数都是属于 W 的特征标的线性组合的一个弱极限。
\end{enumerate}

为证明 \(a\) )，由双极定理，只须证明：若 \(f\) 是 \(\mathrm{L}^{1}(\mathrm{G})\) 中与 U 的元素正交的函数，则 \(f\) 与 W 正交。而此时 Fourier 余变换 \(\overline{\mathcal{F}}(f)\) 在 \(\mathrm{A}(\mathrm{W}) = \mathrm{V}(\mathrm{W}^{\circ})\) 的邻域 U 上为零，故 II，第 252 页命题 2 确实表明 \(f \in \mathrm{W}^{\circ}\)。断言 \(b\) ) 用类似的方式证明，只是把 II，第 252 页命题 2 换成 II，第 257 页定理 2。

